\section{Datos, evaluación y escala como límites del desempeño}
\label{sec:disc_limites}
% Extensión estimada: 2.5 páginas, cinco párrafos (ficha 6.4).
% Fuentes: sec:res_sup_arquitecturas, sec:res_sup_clases, tab:res_supervisados (ensamble), sec:res_melanoma, sec:res_fuga,
% sec:res_comparacion, sec:res_estado_arte, sec:res_gen_nivel1, sec:res_gen_sin_anclaje, sec:res_gen_estado_arte,
% sec:res_costo_clasificadores, sec:res_costo, sec:res_gen_comparacion, tab:marco_generadores, sec:estado_arte_vlm;
% kim2025skinsavvy2, yao2022single, khan2025ensemble, yan2026dermfmzero. Reemplaza las antiguas 6.5 y 6.6.

Si el reporte hereda el desempeño del clasificador, la pregunta siguiente es qué limita ese desempeño. Entre Swin-Tiny, ConvNeXt-Tiny y ViT-B/16 no se encontraron diferencias estadísticamente respaldadas bajo las condiciones evaluadas (Sección~\ref{sec:res_sup_arquitecturas}). Por lo tanto, los resultados no permiten establecer una ventaja clara de una arquitectura sobre las otras, aunque tampoco demuestran que su desempeño sea equivalente. Además, ViT-B/16 no supera a Swin-Tiny con el triple de parámetros, aunque usa la configuración elegida para este último, y el ensamble de los dos mejores supera la media de cada uno en unas 0.02 de F1 macro, una diferencia del orden del desvío entre ejecuciones (Tabla~\ref{tab:res_supervisados}). Una explicación posible es que la disponibilidad y distribución de las imágenes de entrenamiento limiten las mejoras obtenidas mediante arquitecturas de mayor capacidad. No obstante, los experimentos realizados no permiten determinar el peso relativo de ambos factores.

Además de los factores que pueden influir sobre el desempeño real de los clasificadores, las decisiones adoptadas durante la evaluación condicionan las métricas obtenidas y la posibilidad de compararlas con otros trabajos. Estas decisiones mueven las cifras en magnitudes iguales o mayores que las que separan a las arquitecturas. La partición por imagen sobreestima la sensibilidad macro en unas 0.07 (Sección~\ref{sec:res_fuga}),
una cifra del mismo orden que la distancia entre Swin-Tiny y el trabajo de \textcite{kim2025skinsavvy2}, que divide las
imágenes al azar. En melanoma, la regla de decisión pesa más que el modelo: el umbral fijado en validación lleva la
detección de tres a cuatro de cada cinco melanomas sin cambiar el \gls{auc} (Sección~\ref{sec:res_melanoma}), mientras
que los trabajos revisados operan en puntos que van desde la clase más probable hasta una sensibilidad cercana a la
total. Y el rango de la sensibilidad entre ejecuciones supera una décima en las clases minoritarias, de modo que un resultado de una
sola ejecución puede quedar arriba o abajo de otro por la variación entre entrenamientos. Por eso una comparación con la literatura es interpretable cuando se conocen la partición, la regla de decisión y la cantidad de ejecuciones, y por eso este trabajo fija el
umbral en validación, a diferencia de un modelo elegido sobre el mismo conjunto de prueba en que se informa, como el de
\textcite{yao2022single} (Sección~\ref{sec:res_estado_arte}).

La elección de la métrica también puede llevar a conclusiones distintas, especialmente en conjuntos desbalanceados. McNemar compara las imágenes que acierta cada modelo, y el \emph{bootstrap} pareado compara su sensibilidad macro. Según McNemar, DenseNet121 acierta más imágenes que DermLIP, 250 frente a 85 discordantes. Según el \emph{bootstrap}, en cambio, su sensibilidad macro no se distingue de la de DermLIP (Sección~\ref{sec:res_comparacion}). Con clases desbalanceadas, el total de aciertos queda dominado por las clases mayoritarias, y por eso la comparación tiene que aplicarse a una métrica macro.

Una dificultad similar aparece en la evaluación de los reportes generados, aunque las métricas empleadas miden propiedades diferentes. Las medidas de similitud textual permiten comparar los reportes con referencias, mientras que la consistencia diagnóstica identifica la presencia de la categoría verdadera. Ninguna de estas medidas, por sí sola, permite establecer la fidelidad clínica completa del texto. Las métricas basadas en referencias son sensibles a las diferencias léxicas y semánticas entre el reporte generado y el texto de comparación, por lo que sus resultados pueden estar condicionados por el estilo y el contenido de las referencias disponibles: SkinVL, ajustado con las descripciones de \emph{MM-Skin}, obtiene en su propio conjunto de prueba y en una comparación no directa un BLEU-4 más de cien veces
mayor que el de los generadores de este trabajo con un BERTScore parecido (Sección~\ref{sec:res_gen_estado_arte}). La
consistencia diagnóstica, a su vez, favorece enumerar, y por eso LLaVA-Med obtiene sin anclaje la consistencia más alta
nombrando dos clases por reporte (Sección~\ref{sec:res_gen_sin_anclaje}). Conviene, entonces, que una métrica de fidelidad se acompañe de una medida de cuánto se compromete el reporte, como la proporción de reportes que nombran solo la clase verdadera.

Las comparaciones con el estado del arte muestran que el mayor uso de recursos computacionales no garantiza, por sí solo, un mejor desempeño bajo cualquier protocolo. Sin embargo, las diferencias entre trabajos impiden determinar cuánto de la distancia observada puede atribuirse específicamente a la escala. El \emph{pipeline} se ajusta en una sola GPU de 24~GB y su generación entra en una de consumo de 12~GB: Mini-VLM ocupa como máximo 8.6~GB en la A30 y 8.3~GB en la RTX 4070, y en la A30 escribe un reporte en unos 16 segundos, y el clasificador y el \gls{knn} agregan un costo mínimo (Secciones~\ref{sec:res_costo_clasificadores}
y~\ref{sec:res_costo}). Los trabajos revisados, en cambio, ajustan sus generadores en ocho GPU, preentrenan sus modelos
fundacionales en cuatro H100 o redactan con un modelo de 70 mil millones de parámetros consultado por
API~\cite{khan2025ensemble} (Tabla~\ref{tab:marco_generadores} del Capítulo~\ref{cp:marco}). En la clasificación, esa
diferencia no parece decisiva, porque la distancia con los mejores resultados de la literatura es del orden del efecto
de la partición, y entre los modelos evaluados el más grande no mejora a Swin-Tiny. En la representación contrastiva, en cambio, sí parece serlo: DermFM-Zero, preentrenado con tres millones de imágenes en cuatro H100 y alineado después con un millón de pares
imagen--texto, supera a DermLIP en el protocolo de
sus autores con 0.733 frente a 0.561~\cite{yan2026dermfmzero}. Entre los generadores evaluados, por último, el tamaño no ordena la
fidelidad: LLaVA-Med, con casi el doble de parámetros que Mini-VLM, transcribe peor el diagnóstico, y SkinVL, del mismo
tamaño que LLaVA-Med, describe la mitad de las entidades clínicas (Sección~\ref{sec:res_gen_comparacion}).

En las condiciones experimentales evaluadas, la restricción de utilizar una sola GPU con hasta 24~GB de memoria resultó compatible con la ejecución del \emph{pipeline} y con la consistencia diagnóstica obtenida. Este resultado se encuentra asociado a la configuración particular del clasificador y del generador seleccionados. Lo que la restricción deja fuera son dos
cosas que este trabajo no evaluó: los generadores de mayor tamaño, cuya descripción podría ser distinta, y el
preentrenamiento contrastivo a gran escala, que es la vía por la que la literatura mejora la representación que aquí
usa el \gls{knn}.

En conjunto, los resultados sugieren que el desempeño observado está condicionado por una combinación de factores. La disponibilidad y distribución de los datos pueden limitar la capacidad de aprendizaje de los clasificadores, las decisiones de evaluación influyen sobre las métricas obtenidas y las restricciones computacionales determinan qué modelos y procedimientos pueden estudiarse. Aunque los experimentos no permiten aislar la contribución de cada factor, muestran la importancia de considerar estas condiciones al interpretar los resultados y compararlos con los de la literatura.
